\subsection{Curvature pairing and an integrability estimate}

We have proved that $\Theta$ vanishes on every component of
$\mathcal E=\operatorname{div}(s)$. Put
$\mu=i^{n^2}\Theta\wedge\overline\Theta$.
The resulting identity $[\mathcal E]\wedge\mu=0$ will control
singular metrics on $K_{P/C}$, and hence the poles of the normalized
kernel field $v$.

For a semipositive weight $T$ on $K_{P/C}$, the curvature pairing
with $\mu$ has the same integral as $[\mathcal E]\wedge\mu$ and
therefore vanishes. When $T\geq\log|s|^2$, applying this observation
to truncations shows that $\psi=T-\log|s|^2$ is constant along
almost every local kernel leaf. This permits the weighted use of
\eqref{cc:pairing} and gives an integrability exponent strictly
greater than one. We prove the estimate for arbitrary $T$ before
constructing the weight supplied by pluricanonical section spaces.

We use weights for metrics written as \(e^{-\varphi}\), and the convention
\[
 \ddc=\frac{i}{2\pi}\partial\bar\partial,
 \qquad \ddc\log|z|^2=[z=0].
\]
In particular, if \(\varphi_H\) is a local weight for the metric on
\(H\), then \(\ddc\varphi_H=h\).  Expressions such as
\(\log|s|^2\) below denote local weights in holomorphic frames; the
difference of two weights on the same bundle is a globally defined
function wherever both are finite.

\begin{lemma}\label{ca:integrability}
Let \(T\) be a semipositive singular weight on \(K_{P/C}\) such that
\(T\geq\log|s|^2\), and put
\(\psi=T-\log|s|^2\geq0\) off \(\mathcal E\).
There exists \(\delta>0\) such that
\begin{equation}\label{ca:global-integral}
 \int_P e^{(1+\delta)\psi}
       i^{n^2}h\wedge\Theta\wedge\overline\Theta<\infty.
\end{equation}
\end{lemma}

\begin{proof}
The form \(\mu\) is smooth, closed, and positive of type \((n,n)\).
Lemma~\ref{cc:zeros} implies that
\([\mathcal E]\wedge\mu=0\): the integral over each component can be
computed on a resolution of that component, where the pullback of
\(\Theta\) vanishes.  Since the curvature current of
\(\log|s|^2\) is \([\mathcal E]\), cohomology and closedness give
\[
 \int_P (\ddc T)\wedge\mu
    =\int_P[\mathcal E]\wedge\mu=0.
\]
Here and subsequently \(\ddc T\) denotes the globally defined curvature
current of the bundle weight.  The measure on the left is positive,
so it vanishes.  For every real \(b\), the weight
\[
 T_b=\max\{T,\log|s|^2+b\}
\]
is another semipositive weight on the same bundle.  The same argument
shows that \((\ddc T_b)\wedge\mu=0\).

Delete \(\mathcal E\), the zeros of \(\Theta\), and the bad-fibration
locus.  In a holomorphic flow box for the kernel of \(\Theta\), use
coordinates \((w_1,\ldots,w_n,\tau)\), where the leaves vary only in
\(\tau\).  Closedness gives
\[
 \Theta=a(w)\,dw_1\wedge\cdots\wedge dw_n,
 \qquad a(w)\ne0.
\]
The weight \(\log|s|^2\) is pluriharmonic here.  Thus wedging the
curvature of \(T\), or of \(T_b\), with \(\mu\) selects the
\(\tau\)-Laplacian of \(\psi\), or of \(\max\{\psi,b\}\), against
a nonvanishing transverse volume.  Slicing these zero positive
measures shows that both functions are harmonic on almost every
local leaf.  Take the intersection of these full-measure sets for
rational \(b\).  If a harmonic function is nonconstant, there is a
small leaf neighborhood on which its differential is nonzero and
whose image contains an open real interval.  Choose a rational
\(b\) in that interval.  Taking the maximum with \(b\) produces a
nonzero positive Laplacian on the level crossing in that
neighborhood.  Hence \(\psi\) is constant on almost every local
leaf.

In these coordinates $\overline\Theta$ contains all the transverse
differentials $d\bar w_j$. Thus only a derivative in the leaf
direction can contribute to
$\bar\partial(\chi(\psi)\overline\Theta)$, and leafwise constancy
gives, for every bounded continuous function \(\chi\),
\[
 \bar\partial\bigl(\chi(\psi)\overline\Theta\bigr)=0
\]
off the deleted analytic set, in the sense of distributions.  This
identity extends across that set.  Indeed, the coefficients of this
current are locally bounded.  Across a smooth analytic stratum of
complex codimension \(q\geq1\), a cutoff at distance \(\epsilon\)
has derivative \(O(\epsilon^{-1})\) on a tube of volume
\(O(\epsilon^{2q})\), so its contribution tends to zero.  Applying
this successively to the smooth strata, and choosing the tubes about
the remaining lower-dimensional strata sufficiently small, proves
the distributional extension.  This argument uses boundedness of
\(\chi\); it makes no differentiability assumption on \(\psi\).

We may now pair \eqref{cc:pairing} with
\(\chi(\psi)\overline\Theta\).  Choose the representative \(\ell\)
there to be the pullback of a nonnegative smooth area form supported
in relatively compact good base charts.  This does not change its
cohomology class.  On those charts, after shrinking within the good
locus, fiberwise rank one gives
\[
 \alpha|_{D_t}=a(t)s|_{D_t}
\]
with \(a(t)\) bounded on the support in question.  Dolbeault
integration by parts against the closed current just constructed
therefore gives
\begin{equation}\label{ca:paired-truncation}
 \int_P \chi(\psi)\,h\wedge\Theta\wedge\overline\Theta
 =\int_P\chi(\psi)\,\ell\wedge\alpha\wedge\overline\Theta.
\end{equation}
On a good fiber write \(\omega=s|_{D_t}\).  In local frames the
absolute value of the right-hand integrand for
\(\chi(x)=\min\{M,e^{(1+\delta)x}\}\) is bounded by a fixed multiple of
\[
 |\omega|^2 e^{(1+\delta)(T-\log|\omega|^2)}
   =e^{(1+\delta)T}|\omega|^{-2\delta}.
\]
Here \(|\omega|^2\) denotes the canonical volume density.
The weights \(T\) have local upper bounds.  The zero divisor of
\(\omega\) is relatively simple normal crossing on this compact
set, so the last expression is integrable if \(\delta>0\) is small
enough that \(\delta e_j<1\) for its finitely many multiplicities.
The bound is independent of \(M\).  Multiply
\eqref{ca:paired-truncation} by \(i^{n^2}\), take absolute values on
the right, and let \(M\to\infty\).  Monotone convergence proves
\eqref{ca:global-integral}.
\end{proof}

\subsection{Fixed orders and Bergman weights}

Calculate the following orders first on the geometric generic fiber.
The section-space and vanishing-filtration comparisons below will
give the same values on every simultaneous very general fiber.
For a prime divisor \(S\) on such a fiber \(D\) and positive rational
\(\epsilon\), define its actual asymptotic fixed order by
\[
 a_S(\epsilon)=
 \inf_{\substack{m>0\\m\epsilon\in\Z}}
 \frac{1}{m}\min_{0\ne U\in
 H^0(D,mK_D+m\epsilon H|_D)}\ord_S(U),
\]
where degrees with zero section space are omitted.  These spaces
are nonzero in sufficiently divisible degree, since \(K_D\sim E\)
and \(\epsilon H|_D\) is ample.  Multiplication by an ample
section not vanishing generically on \(S\) shows that
\(a_S(\epsilon)\) is nondecreasing as \(\epsilon\) decreases.
In making this comparison, first pass to common divisible degrees
and raise any chosen section to a power.  Thus the limit
\[
 \sigma_S(K_D)=\lim_{\epsilon\downarrow0}a_S(\epsilon)
\]
exists. We need only this definition and its elementary
monotonicity.

Our objective is to prove that $\sigma_S(K_D)>0$ whenever
$S$ comes from a horizontal pole divisor of $v$. A zero value
would give section degrees $k$ and positive ample twists $r_k$
for which both $r_k/k$ and the normalized minimal vanishing
order $j_k/k$ tend to zero. The next two lemmas associate to
such degrees a semipositive weight with zero generic divisorial
coefficient along $S$. The pole-density estimate will then show
that this weight cannot satisfy Lemma~\ref{ca:integrability}
at a pole of $v$.

The weight must be constructed on $P$ before we select the fiber
on which its global integral is finite. We therefore use relative
section spaces on the fixed model of Lemma~\ref{cc:preparation}.
The weight will work on every fiber satisfying the corresponding
base-change conditions. Only afterwards, in
Proposition~\ref{ca:poles}, will we make the final fiber choice.

For any good fiber \(D\) under consideration, write
\[
 \omega=s|_D,\qquad E=\mathcal E|_D=\sum_j e_jD_j,
 \qquad \varphi_0=\log|\omega|^2.
\]
The geometric divisor components are already defined on the fixed
model of Lemma~\ref{cc:preparation}.  Choose the integer degrees and
relative section spaces on this model; no individual sequence of
generic sections needs to be spread.  For each degree, cohomology and
base change also apply to the section
spaces with prescribed vanishing along these components.  Excluding
the resulting countable union of proper closed subsets fixes all
section dimensions and minimal vanishing orders simultaneously.

\begin{lemma}\label{ca:bergman}
Choose integers \(k\to\infty\) and \(r_k>0\) with
\(r_k/k\to0\), such that
\(H^0(D,kK_D+r_kH|_D)\ne0\) on the generic fiber.  There is a
semipositive weight \(T\) on \(K_{P/C}\), depending only on these
degrees and the fixed family and metrics, with \(T\geq\log|s|^2\).
On every good fiber \(D\) satisfying the simultaneous base-change
conditions for these unfiltered section spaces, for every sequence
\(U_k\in H^0(D,kK_D+r_kH|_D)\) normalized by
\begin{equation}\label{ca:normalization}
 \int_D |U_k|^{2/k}e^{-r_k\varphi_H/k}=1,
\end{equation}
the weights and functions
\begin{equation}\label{ca:section-weights}
 \varphi_k=\frac{\log|U_k|^2-r_k\varphi_H}{k},
 \qquad x_k=\varphi_k-\varphi_0,
 \qquad \psi=T|_D-\varphi_0
\end{equation}
satisfy uniform local upper bounds for \(\varphi_k\) and
\(\limsup_k x_k\leq\psi\) off \(E\).  The integral in
\eqref{ca:normalization} is understood in canonical coordinate
densities, and is independent of those coordinates and bundle
frames.
\end{lemma}

\begin{proof}
Apply the relative \(k\)-Bergman construction to
\(\pi:P\to C\) and the twisting bundle \(r_kH\) with its smooth
positive metric.  The source and base are smooth, the morphism is
projective and surjective with connected fibers, and a nonzero
fiber section exists generically.  Generic base change supplies a
nonempty open on which all such sections extend locally.  The
smooth twisting metric makes every fiber section \(2/k\)-integrable.
Thus Theorem~4.2.7 of \cite{PaunTakayama18}, in the cited arXiv version, supplies
a semipositive weight \(\beta_k\) on \(kK_{P/C}+r_kH\).  Over the
smooth extendibility locus, its exponential is the extremum of
\(|U|^2\) among sections normalized as in
\eqref{ca:normalization}.  No extremal interpretation on exceptional
fibers is needed.

Set
\[
 \tau_k=\frac{\beta_k-r_k\varphi_H}{k}.
\]
These are weights on \(K_{P/C}\), and
\(\ddc\tau_k\geq-(r_k/k)h\).  On a fixed smaller submersion chart,
the mean inequality for the subharmonic function given by the
absolute value of a section coefficient to the power \(2/k\),
together with its normalized integral and the bounded factor
\(e^{-r_k\varphi_H/k}\), gives an upper bound independent of \(k\).
Taking the extremum gives that bound for \(\tau_k\).  The bound
persists across degree-dependent exceptional base values by
plurisubharmonic extension.  Equivalently, the ambient estimate in
Remark~4.3.1 of \cite{PaunTakayama18} gives precisely this uniformity when the
divided twisting weights are bounded above.

For completeness, local boundedness on one fixed nonempty open also
globalizes as follows.  Subtract a fixed smooth reference weight on
\(K_{P/C}\).  The resulting globally defined potentials have a
uniform lower bound for their complex Hessians.  If their suprema
were unbounded, subtracting those suprema and using compactness of
sup-normalized quasi-plurisubharmonic potentials would give a
subsequence converging in \(L^1\) to a finite quasi-plurisubharmonic
function.  On the fixed open, however, the normalized functions
would tend uniformly to minus infinity.  This is impossible.
Hence the weights \(\tau_k\) are uniformly bounded above in local
smooth frames on \(P\).

These constructions use the relative section spaces, without selecting
a final fiber or an individual section.  Take their regularized upper
limit on $P$ and then its maximum with the divisor weight:
\begin{equation}\label{ca:limiting-weight}
 T=\max\left\{\left(\limsup_{k\to\infty}\tau_k\right)^*,
                         \log|s|^2\right\}.
\end{equation}
The curvature losses tend to zero, so this is semipositive; if the
first term is identically minus infinity, the second term still
defines the weight.  Its restriction to a good fiber is well
defined because it dominates \(\varphi_0\).  Every normalized
section is bounded above by the corresponding Bergman extremum.
Consequently \(\varphi_k\leq\tau_k|_D\), which proves all the
assertions.
\end{proof}

The weight $T$ now bounds the upper limit of the normalized section
weights in the chosen degrees. We next compare its divisorial coefficient with the minimal
vanishing orders of those complete section spaces. This comparison
holds on every fiber satisfying the stated generic conditions;
the final choice of a fiber on which the integral in
\eqref{ca:global-integral} is finite is not needed for the next lemma.

\begin{lemma}\label{ca:jet}
Let \(D\) satisfy the simultaneous generic conditions above, and fix
a component \(S=D_i\) of \(E\).  Suppose the degrees chosen in
Lemma~\ref{ca:bergman} have minimal actual
vanishing orders
\[
 j_k=\min\{\ord_S(U):0\ne U\in
                  H^0(D,kK_D+r_kH|_D)\}
\]
satisfying \(j_k/k\to0\).  The weight \(T\) constructed in
\eqref{ca:limiting-weight} has zero generic divisorial Lelong
coefficient along \(S\) after restriction to \(D\).
\end{lemma}

If $T|_D$ had positive divisorial coefficient along $S$, an adjoint
estimate would produce a section with the same leading coefficient
and a strictly smaller $L^{2/k}$ integral. We prove both the
preservation of that coefficient and a gain proportional to $1/k$;
the dependence on $k$ matters because the exponent $2/k$ tends to zero.

\begin{proof}
Suppose that coefficient is \(\nu>0\), with Lelong coefficients
normalized by \(\log|z|^2\).  For each \(k\), choose a nonzero
leading section in the image of
\[
 H^0(D,kK_D+r_kH|_D-j_kS)
 \longrightarrow
 H^0\bigl(S,(kK_D+r_kH|_D-j_kS)|_S\bigr).
\]
Among sections with this prescribed image, minimize the
\(2/k\)-integral in \eqref{ca:normalization}.  This minimum exists:
the constraint is a nonempty closed affine subset of a
finite-dimensional vector space, and the integral is continuous,
positive away from zero, and homogeneous of positive degree.
It is therefore coercive.  Denote a minimizer by \(U_k\), and
rescale both it and the prescribed image so that
\eqref{ca:normalization} holds.  All lower coefficients along \(S\)
are zero by the definition of \(j_k\).

Choose \(c>0\) so large that, for all sufficiently large \(k\),
\begin{equation}\label{ca:coefficient-choice}
 c\nu-(1+c)\frac{j_k}{k}>1.
\end{equation}
Choose \(p_0>0\) sufficiently small that
\begin{equation}\label{ca:fiber-integral}
 \int_D |\omega|^2e^{(1+p_0)\psi}<\infty.
\end{equation}
This follows already from the local upper bound for \(T|_D\): the
integrand is bounded by a constant times \(|\omega|^{-2p_0}\),
which is integrable when \(p_0e_j<1\) for every component of the
simple normal crossing divisor \(E\).

Fix a smooth convex function \(g:\R\to\R_{\geq0}\) with
\(0\leq g'\leq1\), equal to \(0\) sufficiently far to the left
and to its argument sufficiently far to the right, and with
\(g''>0\) on \([-1,1]\).  Choose a smooth function
\(0\leq\theta\leq1\) that is zero on \((-\infty,-1]\), one
on \([1,\infty)\), and whose derivative is supported in
\([-1,1]\).  For \(b>0\), put
\begin{equation}\label{ca:adjoint-weight}
 \Phi=(k-1-c-p_0)\varphi_k+cT|_D
          +p_0\bigl(\varphi_0+g(x_k-b)\bigr)+r_k\varphi_H.
\end{equation}
This is a weight on the integral bundle
\((k-1)K_D+r_kH|_D\): the coefficients of its canonical weights
sum to \(k-1\), and \(x_k\) is a global function off the divisors.
We henceforth take \(k>1+c+p_0\).  The constants $c,p_0$ are fixed.
For each fixed $k$ and cutoff $b$, we first remove the regularization
used below.  We then estimate the upper limit as $k\to\infty$ with
$b$ fixed, and finally let $b\to\infty$.  The contradiction will use
one sufficiently large $b$ followed by sufficiently large $k$.

\smallskip\noindent\textbf{The curvature and the adjoint estimate.}\enspace
Off \(E\) and the zero divisor of \(U_k\), one has
\(\ddc\varphi_k=-(r_k/k)h\), \(\ddc\varphi_0=0\), and
\(\ddc x_k=-(r_k/k)h\).  Direct differentiation of
\eqref{ca:adjoint-weight} gives
\begin{align}
 \ddc\Phi={}&\frac{r_k}{k}(1+c+p_0-p_0g'(x_k-b))h
       +c\ddc(T|_D) \notag\\
 &+p_0g''(x_k-b)\frac{i}{2\pi}
               \partial x_k\wedge\bar\partial x_k.
                                                   \label{ca:curvature}
\end{align}
In particular it dominates \((1+c)r_kh/k\) and retains the displayed
gradient term.  The resulting adjoint estimate produces
\(V_k\in H^0(D,kK_D+r_kH|_D)\) such that, on the complement of
the divisors,
\[
 u_k=\theta(x_k-b)U_k-V_k
\]
satisfies
\begin{equation}\label{ca:l2-estimate}
 \int_D\left|\frac{u_k}{U_k}\right|^2w_k
 \leq C_{p_0}\int_{\{|x_k-b|\leq1\}}w_k,
 \qquad
 w_k=|\omega|^2
 e^{(1+c+p_0)x_k-c\psi-p_0g(x_k-b)}.
\end{equation}
The constant is independent of \(b\) and \(k\): the retained gradient
term controls the datum on the cutoff shell, even as the ample
coefficient tends to zero. We give the
singular-weight justification, including the choice of complete
metric, before using this inequality.

Subtract a fixed smooth reference metric from \(T|_D\) and apply
Theorem~14.12(a),(c) and Corollary~14.13(c) of \cite{Demailly12}.  We obtain
weights \(T_\ell\) converging almost everywhere to \(T|_D\),
uniformly bounded above locally, smooth outside an analytic set
\(Z_\ell\), and satisfying
\(\ddc T_\ell\geq-\epsilon_\ell h\), where
\(\epsilon_\ell\to0\).  The cited construction gives pointwise
approximation above the original potential.  In any event, taking
a smooth regularized maximum with \(\varphi_0\), of width tending
to zero, ensures \(T_\ell\geq\varphi_0\).  The regularized
maximum is invariant under simultaneous translation of its
arguments, so it gives a weight on the same bundle.  Its
monotonicity and convexity preserve the lower curvature bound.
It is smooth off \(Z_\ell\cup E\); preservation of analytic
singularities of the maximum itself is not required.

Keep \(k,b\) fixed and replace \(T|_D\) by \(T_\ell\) in
\eqref{ca:adjoint-weight}, obtaining \(\Phi_\ell\).
For large \(\ell\), the loss \(c\epsilon_\ell h\) is less than
half the ample term in \eqref{ca:curvature}.  On
\[
 D\setminus\bigl(Z_\ell\cup E\cup\operatorname{div}(U_k)\bigr)
\]
the weight is smooth with positive curvature and the same gradient
lower bound.  This complement admits a complete K\"ahler metric.
Indeed its deleted analytic set is the zero set of a section of a
Hermitian vector bundle: twist its coherent ideal by a sufficiently
ample bundle and use finitely many generators.  Lemma~12.9 of
\cite{Demailly12} applies, because \(D\) is compact K\"ahler and
the curvature of that smooth bundle has a uniform upper bound.

Write \(A_\ell=i\partial\bar\partial\Phi_\ell\) for the full
positive curvature form on the complement.  Add \(A_\ell\) to a
complete K\"ahler metric there, obtaining a complete metric
\(\widehat h_\ell\geq A_\ell\).  This particular choice also
verifies the ordinary \(L^2\) hypothesis for the datum
\[
 f_k=\theta'(x_k-b)\bar\partial x_k\,U_k.
\]
To see both this and the uniform estimate explicitly, let
\(G\) be the matrix of an auxiliary K\"ahler metric and \(A\)
the positive curvature matrix.  For a bundle-valued \((n,1)\)-form
with local coefficient \(U\eta\), matrix cancellation gives
\begin{equation}\label{ca:metric-cancellation}
 \bigl\langle[A,\Lambda_G]^{-1}(U\eta),U\eta\bigr\rangle_G
       e^{-\Phi}\,dV_G
 = |U|^2e^{-\Phi}\,\eta^*A^{-1}\eta\,d\lambda,
\end{equation}
up to the fixed convention for the canonical volume form.
The determinants from the top holomorphic degree and the volume
cancel, and so do the metric factors in the antiholomorphic degree
and in the curvature operator.  The solution norm in bidegree
\((n,0)\) is likewise independent of \(G\).  When \(G\geq A\),
the ordinary datum norm is bounded by the right side of
\eqref{ca:metric-cancellation}.

The gradient term in \eqref{ca:curvature} implies, for
\(\eta=\bar\partial x_k\), the matrix bound
\[
 \eta^*A_\ell^{-1}\eta
       \leq\frac{1}{p_0g''(x_k-b)}.
\]
Since \(\theta'\) is supported where \(g''\) has a positive
minimum, the inverse-curvature density of \(f_k\) is at most a
constant times the shell density.  That density is integrable.
The datum is \(\bar\partial\)-closed, being the derivative of
\(\theta(x_k-b)U_k\) on the complement.  The complete K\"ahler
estimate, Theorem~5.1 of \cite{Demailly12}, consequently gives
\(\bar\partial u_{k,\ell}=f_k\) with the desired bound for
\(\Phi_\ell\).

Here the local adjoint density has precisely the required
exponents:
\begin{align*}
 |U_k|^2e^{-\Phi_\ell}
 &=\exp\bigl((1+c+p_0)\varphi_k-cT_\ell
                 -p_0(\varphi_0+g(x_k-b))\bigr)\,d\lambda\\
 &=|\omega|^2
       e^{(1+c+p_0)x_k-c(T_\ell-\varphi_0)-p_0g(x_k-b)}.
\end{align*}
On the fixed shell, it is bounded above by
\(e^{(1+c+p_0)(b+1)}|\omega|^2\), because
\(T_\ell-\varphi_0\geq0\).  Moreover \(\Phi_\ell\) is bounded
above locally uniformly in \(\ell\), using the positive coefficient
\(k-1-c-p_0\) and
\[
 \varphi_0+g(\varphi_k-\varphi_0-b)
       \leq\max\{\varphi_0,\varphi_k-b\}+C_g.
\]
The weighted solution estimate therefore bounds the ordinary local
\(L^2\) norms of \(u_{k,\ell}\).  The holomorphic section
\(V_{k,\ell}=\theta(x_k-b)U_k-u_{k,\ell}\) extends across its
deleted analytic set by the distributional \(L^2\) removable-singularity
lemma \cite[Lemma~12.10]{Demailly12}. In a local bundle frame the
right side of the distributional $\bar\partial$ equation is zero,
so the extended coefficients have holomorphic representatives.
The extensions lie in the fixed finite-dimensional space
\(H^0(D,kK_D+r_kH|_D)\), with bounded ordinary norm.  Take a
subsequence converging as global sections.  Fatou's lemma for the
left side and dominated convergence on the fixed shell give
\eqref{ca:l2-estimate}.  For this passage one may discard the
countable union of the sets \(Z_\ell\); no nesting of those sets is
needed.

\smallskip\noindent\textbf{The shell becomes small.}\enspace
Put \(B=1+c+p_0\).  For fixed \(b\), the shell densities in
\eqref{ca:l2-estimate} are bounded by
\(e^{B(b+1)}|\omega|^2\).  At a point off \(E\), a subsequence
contributing to their limsup has a further subsequence with
\(x_k\to x_*\in[b-1,b+1]\).  Lemma~\ref{ca:bergman} gives
\(x_*\leq\psi\); thus \(\psi\geq b-1\) and
\[
 Bx_*-c\psi-p_0g(x_*-b)\leq(1+p_0)\psi.
\]
Reverse Fatou therefore yields
\begin{equation}\label{ca:shell-tail}
 \limsup_{k\to\infty}\int_{\{|x_k-b|\leq1\}}w_k
 \leq\int_{\{\psi\geq b-1\}}|\omega|^2e^{(1+p_0)\psi}.
\end{equation}
By \eqref{ca:fiber-integral}, the right side tends to zero as
\(b\to\infty\).

\smallskip\noindent\textbf{The prescribed leading section is preserved.}\enspace
Near a general point of \(S=(z=0)\), write
\(U_k=z^{j_k}a_k\) and \(\omega=z^{e_i}a_0\), with the
coefficients nonvanishing in the chosen frames.  Since
\(j_k/k<e_i\), the function \(x_k\) tends to plus infinity as
\(z\to0\).  Thus \(\theta(x_k-b)=1\) and
\(g(x_k-b)=x_k-b\) sufficiently near that point.
The divisorial decomposition of the positive curvature current of
\(T|_D\) gives
\(T|_D\leq\nu\log|z|^2+O(1)\): subtracting
\(\nu[S]\) leaves a positive current, whose local potential is
bounded above.  See \cite[Chapter~III, Proposition~8.16]{DemaillyCADG}.
Consequently
\begin{equation}\label{ca:jet-weight-order}
 \Phi\leq\left((k-1-c)\frac{j_k}{k}+c\nu\right)
                      \log|z|^2+O(1).
\end{equation}
The coefficient exceeds \(j_k+1\) by
\eqref{ca:coefficient-choice}.  In this neighborhood
\(u_k=U_k-V_k\) is holomorphic.  If its vanishing order were at
most \(j_k\), the transverse integral against \(e^{-\Phi}\)
would diverge, contradicting \eqref{ca:l2-estimate}.  Hence
\(U_k-V_k\) vanishes along \(S\) to order at least \(j_k+1\).
This divisorial statement says exactly that \(V_k\) has the
prescribed leading section and all the prescribed zero lower
coefficients.

\smallskip\noindent\textbf{A gain of order \(1/k\).}\enspace
Let \(\rho_k\,dV\) be the normalized density in
\eqref{ca:normalization} relative to a fixed smooth volume on
\(D\), and put \(q_k=|u_k/U_k|\) off the zero divisor.
The functions \(\rho_k\) have a uniform upper bound by
Lemma~\ref{ca:bergman}. In local canonical frames the density is
\[
 w_k=\exp\bigl(B\varphi_k-cT|_D
             -p_0(\varphi_0+g(x_k-b))\bigr)\,d\lambda.
\]
The local upper bounds for \(T|_D\), and the earlier bound
$\varphi_0+g(x_k-b)\leq\max\{\varphi_0,\varphi_k-b\}+C_g$,
therefore give, on a finite cover by smooth frames,
\begin{equation}\label{ca:density-lower-bound}
 w_k\geq C^{-1}\rho_k^B\,dV
\end{equation}
with \(C\) independent of \(k\) and of \(b\) bounded below.
Fix \(0<\eta<\min\{1,2/B\}\).  H\"older's inequality gives
\begin{align}
 \int_D q_k^\eta\rho_k\,dV
 &\leq
 \left(\int_D q_k^2\rho_k^B\,dV\right)^{\eta/2}
 \left(\int_D
       \rho_k^{(2-B\eta)/(2-\eta)}\,dV\right)^{1-\eta/2}
                                                        \label{ca:holder}\\
 &\leq C_\eta
       \left(\int_D q_k^2w_k\right)^{\eta/2}.\notag
\end{align}
The exponent in the second factor is positive, so its uniform
bound follows from compactness and the upper bound for \(\rho_k\).
Equations~\eqref{ca:l2-estimate} and \eqref{ca:shell-tail} show
that this moment can be made arbitrarily small by first fixing
\(b\) sufficiently large and then taking \(k\) sufficiently large.

The set where \(\theta(x_k-b)\ne0\) is contained in
\(\{\varphi_0\leq C-b+1\}\), by the upper bound for
\(\varphi_k\).  Its smooth volume tends to zero with \(b\),
and therefore its \(\rho_k\)-mass tends to zero uniformly in
\(k\).  Markov's inequality applied to
\eqref{ca:holder} now gives a set of \(\rho_k\)-mass at least
\(1/2\) on which \(\theta(x_k-b)=0\) and \(q_k\leq1/2\).
On this set \(|V_k/U_k|=q_k\leq1/2\); everywhere else,
\(|V_k/U_k|\leq1+q_k\).

For \(0<p\leq\eta/2\) and \(q\geq0\),
\begin{equation}\label{ca:small-exponent}
 (1+q)^p\leq1+C_\eta p q^\eta.
\end{equation}
For \(q\leq1\), this follows from concavity.  For \(q\geq1\),
integrate the derivative with respect to the exponent and use
\((1+q)^{\eta/2}\log(1+q)\leq C_\eta q^\eta\).
Taking \(p=2/k\), we therefore obtain
\begin{align}
 \int_D |V_k|^{2/k}e^{-r_k\varphi_H/k}
 &=\int_D |V_k/U_k|^p\rho_k\,dV\notag\\
 &\leq1-\frac12(1-2^{-p})
             +C_\eta p\int_D q_k^\eta\rho_k\,dV<1
                                                        \label{ca:strict-gain}
\end{align}
for the indicated choices of \(b\) and then \(k\).
Indeed \((1-2^{-p})/p\to\log2\), whereas the moment can be
made arbitrarily small independently of large \(k\).
This is a strict improvement of the minimizing integral for a
section with the same prescription, a contradiction.  Therefore
\(\nu=0\).
\end{proof}

\subsection{Excluding poles with zero limiting fixed order}

We now apply the estimates to a horizontal pole of $v$.
The degrees and the limiting weight will be chosen from the
geometric generic section spaces. Only after applying
Lemma~\ref{ca:integrability} to that fixed weight will we choose
a complex fiber and the minimizing sections on it.

\begin{proposition}\label{ca:poles}
Every horizontal pole divisor of \(v\) has positive limiting fixed
order on a simultaneous very general fiber.  In particular, for
one sufficiently small positive rational \(\epsilon\), every such
divisor has positive actual asymptotic fixed order for
\(K_D+\epsilon H|_D\).
\end{proposition}

\begin{proof}
Suppose a horizontal pole divisor induces \(S=D_i\) and has
zero limiting fixed order on the geometric generic fiber.  Choose
rational perturbations tending to zero and degrees in which the
normalized minimal orders approach their infima.  Passing to powers
as necessary gives integers
\(k\to\infty\), \(r_k>0\), with
\[
 \frac{r_k}{k}\longrightarrow0,
 \qquad \frac{j_k}{k}\longrightarrow0,
\]
where \(j_k\) is the minimal actual order of the whole section
space.  Use the corresponding relative section spaces and vanishing
filtrations on the fixed prepared family, and impose the countably many
base-change conditions that preserve these orders.  No further field
extension or spreading of individual sections is required.
Lemma~\ref{ca:bergman} now constructs one global weight $T$ from the
chosen degrees, before any final fiber is selected.

Apply Lemma~\ref{ca:integrability} to this $T$.  Fubini gives a
full-measure set of good fibers on which the integral is finite.
The excluded algebraic conditions form a countable subset of the
complex curve and hence have measure zero.  Choose $D$ satisfying
both requirements.  Lemma~\ref{ca:jet} applies on this same fiber
and proves that $T|_D$ has zero generic divisorial coefficient at $S$.
The normalized minimizing sections in that Lemma are chosen only now;
they do not change $T$.

In a smooth chart over the base, contraction gives
\[
 \Theta=\iota_v(dt\wedge\omega).
\]
The density of \(i^{n^2}h\wedge\Theta\wedge\overline\Theta\)
relative to base area is comparable, by the smooth positive metric,
to \(|v|_h^2|\omega|^2\).  This formula includes the mixed
components of the metric.  At a generic point of a pole of order
at least one, shrink the chart so its leading coefficient is
bounded below.  Since \(\ord_S(\omega)=e_i\), the finite fiber
integral from \eqref{ca:global-integral} would then dominate a
positive constant times
\begin{equation}\label{ca:pole-integral}
 \int |z|^{-2-2\delta e_i}e^{(1+\delta)T|_D}\,d\lambda,
 \qquad S=(z=0).
\end{equation}
We show that this last integral is infinite.

The generic Lelong coefficient is zero.  The analyticity of Lelong
upper-level sets, or equivalently the generic-Lelong-number
statement in \cite[Chapter~III, Lemma~8.15]{DemaillyCADG}, implies
that the point Lelong number is zero at almost every center on a
small smooth patch of \(S\).  For a local plurisubharmonic weight
\(T\), the logarithmic mean characterization
\cite[Chapter~III, Corollary~6.8 and Example~6.9]{DemaillyCADG}
then gives
\[
 \frac{1}{\vol(B(a,r))}\int_{B(a,r)}T\,d\lambda
                   =o(|\log r|)
\]
at almost every such center \(a\).  The characterization for ball
means follows from that for sphere means by radial integration.

Fix \(\varepsilon>0\).  For almost every center, the ball mean
is at least \(\varepsilon\log r\) at all sufficiently small
dyadic radii.  The sets on which this holds beyond one fixed
dyadic index have a union of full divisor measure.
Hence there is a measurable set \(A\) of positive divisor measure
on a compact patch of \(S\), and one common starting index, for
which the lower bound holds uniformly.  In dimension one the
divisor measure here is counting measure.

On \(B(a,r)\), \(|z|\leq Cr\).  Jensen's inequality therefore
gives, for the nonnegative integrand \(F\) in
\eqref{ca:pole-integral},
\[
 \int_{B(a,r)}F\,d\lambda
 \geq C^{-1}r^{2n-2-2\delta e_i+(1+\delta)\varepsilon}
 \qquad(a\in A).
\]
Integrate this inequality in \(a\).  For each ambient point, the
measure of centers \(a\in A\) whose radius-\(r\) ball contains
it is at most \(Cr^{2n-2}\), by projection to the smooth
tangential coordinates of \(S\).  Fubini consequently yields
\[
 \int F\,d\lambda
 \geq C^{-1}r^{-2\delta e_i+(1+\delta)\varepsilon}.
\]
The overlap bound is constant when \(n=1\), so the same formula
holds in that case.  Since \(e_i\geq1\), choose
\(\varepsilon<2\delta e_i/(1+\delta)\) and let the dyadic
radius tend to zero.  The right side diverges, contradicting
\eqref{ca:pole-integral} and the Fubini consequence of
\eqref{ca:global-integral}.  A pole with zero limiting fixed order
is therefore impossible.

There are only finitely many horizontal pole divisors.  For each,
positivity of its limit and monotonicity give a sufficiently small
rational perturbation with positive actual asymptotic fixed order.
Taking one common smaller rational perturbation proves the last
assertion.
\end{proof}
